\chapter{Related Work}

\section{Literature Review}

A. Overview of the Research Scope

Workplace safety monitoring has evolved significantly with advances in \textbf{computer vision, deep learning, and intelligent video analytics}. Traditional safety methods—manual supervision, periodic inspections, or post-incident video review—are increasingly being augmented by automated systems capable of real-time hazard detection. This chapter reviews literature related to the four core components relevant to VisionSafe 360:

\begin{itemize}
\item \textbf{PPE (Personal Protective Equipment) Detection}
\item \textbf{Human Fall Detection Using Temporal Modelling}
\item \textbf{Proximity and Collision Risk Detection via Multi-Object Tracking}
\item \textbf{Posture and Ergonomic Risk Assessment through Pose Estimation}
\end{itemize}

The aim is to understand state-of-the-art progress, identify existing challenges, and clarify the research gaps that VisionSafe 360 is designed to address.

B. Thematic Grouping of Literature

1. PPE Detection Systems

Study 1 — Ali et al. (2020)

\begin{itemize}
\item \textbf{Goal:} Real-time detection of hardhats and safety vests using CCTV video.
\item \textbf{Method:} YOLOv3-based CNN trained on public construction site datasets.
\item \textbf{Findings:} Achieved \textit{88\% mAP} under controlled lighting conditions.
\item \textbf{Strengths/Limitations:} Robust for helmets; accuracy drops with occlusion or low-light scenarios.
\end{itemize}

Study 2 — Zhang et al. (2021)

\begin{itemize}
\item \textbf{Goal:} Scalable PPE compliance monitoring in large construction sites.
\item \textbf{Method:} YOLOv4 combined with DeepSORT tracking to maintain worker identities.
\item \textbf{Findings:} Consistent detection across multiple camera angles; \textit{82\% ID tracking accuracy}.
\item \textbf{Strengths/Limitations:} High computational demands limit deployment on edge devices.
\end{itemize}

Study 3 — López et al. (2022)

\begin{itemize}
\item \textbf{Goal:} Detect PPE violations (helmet, vest, gloves, boots).
\item \textbf{Method:} Lightweight MobileNet-SSD architecture for edge devices.
\item \textbf{Findings:} Good speed (20–25 FPS on Jetson Nano) with \textit{75\% detection accuracy}.
\item \textbf{Strengths/Limitations:} Limited precision for small PPE items like gloves.
\end{itemize}

2. Fall Detection Using Temporal Modelling

Study 4 — Nguyen et al. (2020)

\begin{itemize}
\item \textbf{Goal:} Fall detection in industrial videos.
\item \textbf{Method:} LSTM networks using frames extracted from CCTV.
\item \textbf{Findings:} \textit{92\% accuracy} on staged falls; reduced reliability in crowded scenes.
\item \textbf{Strengths/Limitations:} Temporal modelling is strong, but dependent on clean video input.
\end{itemize}

Study 5 — Han et al. (2021)

\begin{itemize}
\item \textbf{Goal:} Real-time fall classification.
\item \textbf{Method:} 3D CNN + optical flow to capture temporal patterns.
\item \textbf{Findings:} Works well on multiple camera angles; \textit{86\% precision}.
\item \textbf{Strengths/Limitations:} Computationally expensive; not suitable for low-end GPUs.
\end{itemize}

Study 6 — Wang et al. (2022)

\begin{itemize}
\item \textbf{Goal:} Edge-based fall detection for IoT safety systems.
\item \textbf{Method:} Tiny-YOLO + GRU model for sequences.
\item \textbf{Findings:} 17 FPS edge deployment with reasonable accuracy.
\item \textbf{Strengths/Limitations:} Reduced accuracy during occlusion or distance >8 meters.
\end{itemize}

3. Proximity and Collision Risk Detection

Study 7 — Kim et al. (2021)

\begin{itemize}
\item \textbf{Goal:} Worker-vehicle proximity alerts.
\item \textbf{Method:} DeepSORT multi-object tracking + depth estimation.
\item \textbf{Findings:} Real-time distance estimation; \textit{80\% near-miss prediction accuracy}.
\item \textbf{Strengths/Limitations:} Requires camera calibration for accurate distance estimation.
\end{itemize}

Study 8 — Carter et al. (2022)

\begin{itemize}
\item \textbf{Goal:} Predict collision risks between forklifts and workers.
\item \textbf{Method:} YOLO-based object detection + probabilistic motion prediction.
\item \textbf{Findings:} Reduced false alerts by \textit{25\%} compared to fixed-threshold systems.
\item \textbf{Strengths/Limitations:} Complex motion modelling; sensitive to camera shake.
\end{itemize}

4. Posture and Ergonomics Assessment

Study 9 — Ahmed and Patel (2019)

\begin{itemize}
\item \textbf{Goal:} Detect ergonomically harmful postures in warehouses.
\item \textbf{Method:} OpenPose skeleton extraction + angle-based risk scoring.
\item \textbf{Findings:} Effective at detecting bending, twisting, and lifting risks.
\item \textbf{Strengths/Limitations:} Struggles with partial body visibility.
\end{itemize}

Study 10 — Brown et al. (2021)

\begin{itemize}
\item \textbf{Goal:} Automated ergonomic risk assessment (REBA score prediction).
\item \textbf{Method:} CNN + LSTM for joint angle sequences.
\item \textbf{Findings:} \textit{74\% agreement rate} with expert physiologists.
\item \textbf{Strengths/Limitations:} Requires large annotated posture datasets.
\end{itemize}

C. Summary of the Reviewed Studies

Collectively, prior work demonstrates that:

\begin{itemize}
\item \textbf{Object detection} (YOLO, SSD) is highly effective for PPE compliance but limited by occlusion and lighting.
\item \textbf{Temporal networks} (LSTM, GRU, 3D CNNs) improve fall detection but often require high compute resources.
\item \textbf{Multi-object tracking} aids proximity detection but is sensitive to camera calibration.
\item \textbf{Pose estimation} provides ergonomic insights but struggles with camera viewpoint diversity.
\end{itemize}

The reviewed systems show promising results but often target only \textbf{single hazards}, require \textbf{high-end GPUs}, or depend on \textbf{specialised hardware}.

D. Similar Applications and Existing Solutions.

Although several commercial and research-based systems exist for workplace safety, many exhibit critical limitations.

\newcommand{\companylogowidth}{5.2cm}
\newcommand{\companylogoheight}{2.4cm}
\newcommand{\companylogo}[1]{%
\includegraphics[width=\companylogowidth,height=\companylogoheight,keepaspectratio]{#1}%
}

.1Intenseye

\begin{figure}[H]
\centering
\companylogo{images/img_002.png}
\caption{Intenseye Company}
\end{figure}

is a cloud-based platform focused on monitoring PPE compliance. It offers high detection accuracy and a user-friendly dashboard that simplifies safety management. Nevertheless, the system relies on cloud streaming, making it unsuitable in environments with limited network bandwidth or strict privacy requirements.

2. Protex AI

\begin{figure}[H]
\centering
\companylogo{images/img_003.png}
\caption{Protex Company}
\end{figure}

is another commercial solution that tracks safety incidents such as PPE violations and near-miss events. While it delivers enterprise-grade analytics, it does not support posture or ergonomic assessment, and its high cost can be prohibitive for many companies.

3. SmartVid.io

\begin{figure}[H]
\centering
\companylogo{images/img_004.png}
\caption{StarVid.IO Company}
\end{figure}

is an AI-based safety analytics platform designed for construction sites. It offers robust hazard classification capabilities and can detect unsafe acts effectively, providing safety officers with actionable insights. However, it is a paid service and lacks flexibility for local deployment, which limits its accessibility for smaller organisations or projects in regions with limited resources.

4. YOLO-Based Academic Prototypes.

\begin{figure}[H]
\centering
\companylogo{images/img_005.png}
\caption{Ultralytics Company}
\end{figure}

use object detection models for PPE and hazard detection. Modern versions like \textbf{Ultralytics YOLOv5 and YOLOv8} offer fast, accurate real-time detection and edge deployment. However, most academic implementations focus only on detection and lack dashboards, alerts, or multi-hazard integration, limiting practical workplace use.

\textbf{E. Research Gap}

Despite advancements, several key gaps remain:

1. Lack of Unified Multi-Hazard Systems

Most existing systems address \textbf{only one hazard}—PPE, falls, or proximity—but not a combined pipeline.

2. Limited Real-Time Performance on Affordable Hardware

High-end models demand powerful GPUs, making them inaccessible for many workplaces.

3. Inadequate Support for Non-Ideal Conditions

Challenges include:

\begin{itemize}
\item Low light
\item Occlusion
\item Non-frontal camera angles
\item Crowded industrial scenes
\end{itemize}

4. Absence of Open, Local Deployment Solutions

Many systems require cloud processing, which is unsuitable for:

\begin{itemize}
\item Privacy-sensitive environments
\item Poor network stability
\item Cost constraints
\end{itemize}

5. Lack of Integrated Dashboard \& Analytics

Research often focuses on detection models but ignores real-world usability elements such as:

\begin{itemize}
\item Alert logging
\item Incident querying
\item Live dashboards
\item Worker behavior timelines
\end{itemize}

VisionSafe 360 directly addresses these gaps.

F. Contribution of VisionSafe 360

VisionSafe 360 provides the following contributions:

1. Multi-Hazard Real-Time Monitoring System

A unified pipeline that detects:

\begin{itemize}
\item PPE compliance
\item Fall events
\item Worker–vehicle proximity risks
\item Ergonomic posture deviations
\end{itemize}

2. Optimized for Existing CCTV Infrastructure

No special sensors or IR cameras required; works on standard RGB streams (RTSP).

3. Efficiency-Focused Deep Learning Models

YOLO-based detection + lightweight temporal networks to enable near real-time operation (>15 FPS on common GPUs).

4. Practical Dashboard for Safety Officers

Including:

\begin{itemize}
\item Real-time alerts
\item Event logs
\item Video evidence
\item Hazard analytics
\end{itemize}

\textbf{ 5. Scalable and Modular Architecture}

Each module (PPE, fall, tracking, posture) can be improved independently.

\begin{table}[H]
\centering
\small
\renewcommand{\arraystretch}{1.2}
\begin{tabular}{|
>{\RaggedRight\arraybackslash}p{0.18\linewidth}|
>{\RaggedRight\arraybackslash}p{0.21\linewidth}|
>{\RaggedRight\arraybackslash}p{0.24\linewidth}|
>{\RaggedRight\arraybackslash}p{0.24\linewidth}|}
\hline
\textbf{Feature} & \textbf{Manual Monitoring} & \textbf{Commercial AI Platforms} & \textbf{VisionSafe 360 (Proposed)} \\
\hline
Real-Time Detection & No (Reactive) & Yes & Yes \\
\hline
Multi-Hazard Support & High (Human judgment) & Varies (Usually Modular/Paid) & Unified (PPE, Falls, Proximity, Ergo) \\
\hline
Cost & High (Man-hours) & High (Licensing/Subscription) & Low (Open Source/Standard Hardware) \\
\hline
Data Privacy & N/A & Cloud-based (Data leaves site) & On-Premises (Data stays local) \\
\hline
Customisation & Flexible & Low (Proprietary) & High (Modular Architecture) \\
\hline
\end{tabular}
\caption{\textbf{Research gap vs. VisionSafe 360 contributions}}
\end{table}
